\documentclass[12pt]{article}
\usepackage[utf8]{inputenc}
\usepackage[T1]{fontenc}
\usepackage{hyperref}
\title{Buli Language Resources for K-12 Curriculum Development}
\author{Compiled by Claude Sonnet 5 (claude-sonnet-5) for Andr\'as Kornai}
\date{}
\begin{document}
\maketitle

\section{Introduction}
Buli (ISO 639-3 uly; Glottocode buli1260) is a West Chadic (South Bauchi) Afro-Asiatic language of Bauchi State, Nigeria. Speaker counts conflict: ASJP lists 2,000, a CAPRO survey about 4,000, and Caron (2011) fewer than 1,000 and severely endangered. It should not be confused with Buli of Halmahera, Indonesia, nor with Boghom.

\section{Primary Educational Resources}
None found.

\section{Linguistic and Reference Materials}
\begin{itemize}
\item ASJP wordlist entry, Buli (Nigeria): \url{https://asjp.clld.org/languages/BULI_NIGERIA}
\item Glottolog reference for Cosper (1999), \emph{Barawa Lexicon}, a wordlist covering Buli and seven other South Bauchi languages: \url{https://glottolog.org/resource/reference/id/136865}
\item Wikipedia on South Bauchi languages: \url{https://en.wikipedia.org/wiki/South_Bauchi_languages}
\item Glottolog reference record: \url{https://glottolog.org/resource/reference/id/weball:4586}
\end{itemize}

\section{Religious and Community Resources}
None found.

\section{Digital and Audio Resources}
None found.

\section{Curriculum Development Context}
Only a wordlist and classificatory material were found. Buli speakers use Hausa as a trade language. A native-language curriculum would require new documentation first. Searches were limited, and materials may exist that were not located.
\end{document}